% 问题三子问题二候选正文。
% 仅使用 B6/B7 原生 Q_B 条件模型及已登记的 Q3 优化运行；
% 当前稿件待科学评审，不直接修改 paper/main.tex。

\subsection{问题三子问题二：M1 的规模--数据量--质量条件优化}
\label{sec:q3-subproblem2}

\subsubsection{问题分析与变量界定}

子问题二在问题二 M1 的基础上，将质量资源作为第三个决策变量。本文记该变量为
$Q_B$，它是 B6/B7 数据表中的原生 \texttt{Q\_score}，不是问题一得到的
A 侧质量变量 $Q_A$。因此，$Q_B$ 只在 B6/B7 条件模型内部解释为质量分数，
不能与 $Q_A$ 直接相等，也不能把两个量尺合并成一个已验证的质量变量。

M1 的三个决策变量为模型规模 $N$、训练数据量 $D$ 和原生质量分数 $Q_B$。
其中 $N$ 和 $D$ 分别以十亿参数和十亿 tokens 计，$Q_B$ 的观测范围为
$[0.1,1.0]$。$N$ 和 $D$ 的有界范围取自 B1 样本：
\begin{equation}
  N\in[0.070542,11.965825],\qquad
  D\in[0.134,299.893].
  \label{eq:q3-sp2-support}
\end{equation}
这些区间表示当前数据的观测适用范围，不表示对更大模型或更多训练数据的物理定律外推许可。

\subsubsection{M1 目标函数与质量成本}

冻结的 B6 M1 条件损失为
\begin{equation}
  \widehat L_{\mathrm{M1}}(N,D,Q_B)
  =E+A N^{-\alpha}+B D^{-\beta}+G(1-Q_B).
  \label{eq:q3-sp2-m1-loss}
\end{equation}
本轮使用的参数为
\begin{equation}
\begin{aligned}
E&=1.4892660413, & A&=0.5401729883, & B&=1.3167948616,\\
G&=0.3622474350, & \alpha&=0.2790539239, &
\beta&=0.2828029135.
\end{aligned}
\label{eq:q3-sp2-m1-parameters}
\end{equation}

给定上下文长度 $L_{\mathrm{ctx}}$ 时，基础规模成本写为
\begin{equation}
  C_{\mathrm{base}}
  =10^{18}ND\left(6+\eta L_{\mathrm{ctx}}\right),
  \qquad \eta=2\times10^{-4}.
  \label{eq:q3-sp2-base-cost}
\end{equation}
相对于基准质量 $Q_0=0.1$ 的质量增量成本写为
\begin{equation}
  C_Q=10^9D\left[g(Q_B)-g(Q_0)\right].
  \label{eq:q3-sp2-quality-cost}
\end{equation}
本文比较三类成本形状：
\begin{equation}
\begin{aligned}
g_{\mathrm{exp}}(Q)&=10^7\exp(6Q),\\
g_{\mathrm{pow}}(Q)&=5\times10^9Q^4,\\
g_{\mathrm{log}}(Q)&=2\times10^9\log(1+10Q).
\end{aligned}
\label{eq:q3-sp2-cost-families}
\end{equation}
三类 $g$ 用于敏感性比较，不被解释为已经由真实训练账单验证的成本定律。
在预算 $C$ 下，优化问题为
\begin{equation}
\begin{aligned}
\min_{N,D,Q_B}\quad&
\widehat L_{\mathrm{M1}}(N,D,Q_B),\\
\mathrm{s.t.}\quad&
C_{\mathrm{base}}(N,D;L_{\mathrm{ctx}})
+C_Q(D,Q_B)\le C,\\
&N\in[0.070542,11.965825],\\
&D\in[0.134,299.893],\\
&Q_B\in[0.1,1.0].
\end{aligned}
\label{eq:q3-sp2-opt}
\end{equation}

质量成本占比定义为
\begin{equation}
  \rho_Q=\frac{C_Q}{C_{\mathrm{base}}+C_Q}.
  \label{eq:q3-sp2-quality-share}
\end{equation}
$\rho_Q$ 是当前成本假设下的模型内比例，不是现实工程账单中的实测占比。

\subsubsection{求解方法}

为保持正值约束，$N$ 和 $D$ 采用对数参数化，优化变量为
$(\log N,\log D,Q_B)$。预算约束按预算值缩放后送入 SLSQP，以避免
$10^{19}$--$10^{24}$ 量级的原始约束造成有限差分尺度失衡。

每个情景使用六个确定性初始点，覆盖低值、几何中点、高值以及不同的
$Q_B$ 初始值。多起点的目的在于检查局部求解是否受初始点影响，并比较
可行性、收敛状态、目标函数值和边界状态。最终优先保留优化器报告收敛且预算可行的解。
当前不采用遗传算法或粒子群算法，原因是决策维度低、目标函数连续光滑、约束边界明确，
解析结构与多起点局部优化更便于逐情景审计。

求解后将结果分为三类：$Q_B=0.1$ 或 $Q_B=1.0$ 的端点状态；
$0.1<Q_B<1.0$ 且 $N,D$ 未触及边界的内点状态；以及 $N$、$D$ 或 $Q_B$
触及观测范围边界的边界状态。

\subsubsection{\texorpdfstring{原生 $Q_B$ 情景结果}{原生 Q\_B 情景结果}}

本轮共生成 $3\times5\times3=45$ 个原生 $Q_B$ 情景，覆盖三类质量成本、
五档上下文
\[
L_{\mathrm{ctx}}\in\{2048,4096,8192,32768,131072\}
\]
和三档预算
\[
C\in\{10^{19},10^{22},10^{24}\}\ \mathrm{FLOPs}.
\]
每行结果报告 $Q_B^*$、$N^*$、$D^*$、$L^*$、质量成本占比、预算活跃状态、
边界状态和多起点收敛信息。

按上下文取平均后，低预算 $C=10^{19}$ 下，指数型、幂函数型和对数渐进型
的最优质量分别约为 $0.7243$、$0.6103$ 和 $1.0000$。预算提高到
$10^{22}$ 或 $10^{24}$ 后，当前参数下三类成本通常均将 $Q_B^*$ 推到
$1.0000$。这只说明在当前 M1 和成本假设下质量增量成本相对 $N,D$ 成本较小，
不表示真实训练质量已经饱和，也不表示 $Q_B=1$ 是经验验证的唯一最优质量。

在指数成本、$L_{\mathrm{ctx}}=8192$、$C=10^{22}$ 的代表性情景中，质量成本
约占总成本的 $5.3\%$--$7.4\%$。高预算情景中若同时触及
$N=11.965825$、$D=299.893$ 和 $Q_B=1.0$，则应将其标记为观测范围边界解；
名义预算的剩余部分不能用来证明支持范围外仍有有效收益。

\begin{table}[htbp]
  \centering
  \caption{M1 原生质量条件情景的摘要}
  \label{tab:q3-sp2-summary}
  \small
  \begin{tabular}{lccc}
    \toprule
    质量成本族 & $C=10^{19}$ 的平均 $Q_B^*$ &
    $C=10^{22}$ 的典型状态 & 主要边界 \\
    \midrule
    指数型 & 0.7243 & 接近 $1.0$ & 低预算内点、高预算边界 \\
    幂函数型 & 0.6103 & 接近 $1.0$ & 低预算内点、高预算边界 \\
    对数渐进型 & 1.0000 & $1.0$ & 质量上界较早活跃 \\
    \bottomrule
  \end{tabular}
\end{table}

\subsubsection{上下文成本结构}

由式~\eqref{eq:q3-sp2-base-cost}，基础成本中的注意力项与训练项之比为
\begin{equation}
  \frac{C_{\mathrm{attn}}}{C_{\mathrm{train}}}
  =\frac{\eta L_{\mathrm{ctx}}}{6}
  =\frac{L_{\mathrm{ctx}}}{30000}.
  \label{eq:q3-sp2-context-ratio}
\end{equation}
因此，$L_{\mathrm{ctx}}=30000$ 是当前成本结构中的代数临界点：
低于该值时注意力项小于训练项，高于该值时注意力项大于训练项。
它不是已经由任务效果验证的上下文性能临界点，也不是 $Q_B$ 的经验转折点。

\subsubsection{数值复核与稳健性}

本轮复核包括以下内容：
\begin{enumerate}
  \item 45 行原生 $Q_B$ 输出及参数折叠传播结果均为有限值；
  \item 原生质量情景最大相对预算误差约为 $2.1\times10^{-12}$；
  \item 质量成本稳健性运行的最大相对预算误差约为
  $9.7\times10^{-12}$，参数折叠运行的 M1 最大相对误差约为
  $9.01\times10^{-12}$；
  \item $N,D$ 未超过 B1 观测范围，$Q_B$ 未超过 B6 原生范围；
  \item 六个起点的成功状态、目标值差异和边界标志均被保留；
  \item 原生情景表 SHA256 为
  {\scriptsize\texttt{77171f327c40d0bbd2559665b5aa7b6fdc945b5052301dbd1c73296289330378}}；
  \item $Q_0\in\{0.1,0.3,0.5\}$ 的扫描显示，低预算下指数型和幂函数型
  对 $Q_0$ 有可见敏感性，对数渐进型在当前参数下系统性地推向质量上界。
\end{enumerate}

\subsubsection{结果讨论与小结}

在 B1 的 $N,D$ 观测范围、B6/B7 原生 $Q_B$ 范围以及给定预算、上下文和
质量成本假设下，可以报告质量--规模--数据量的条件替代关系、成本族造成的
资源配置差异、上下文改变后的训练--注意力成本转移，以及参数折叠和边界检查
支持的稳健范围。

不能据此声称 A--B 联合 Loss 已经验证，不能声称 $Q_A$ 与 $Q_B$ 已经建立稳定
实证关系，不能声称 $Q_B=1$ 是真实训练中的唯一最优质量，也不能把支持范围外的
结果解释为现实中的全局最优。以上结论均属于冻结 M1 参数和成本假设下的条件情景。

本子问题的证据运行包括：
\begin{itemize}
  \item \texttt{q3-q-conditional-20260925-r01}；
  \item \texttt{q3-q-robustness-20260925-r01}；
  \item \texttt{q3-optimization-robustness-20260926-r01}；
  \item \texttt{q3-answer-package-20260926-r01}。
\end{itemize}
